% ------------------------------------------------------------------------------
% ABSTRACT
% ------------------------------------------------------------------------------
\chapter*{Abstract}
\addcontentsline{toc}{chapter}{Abstract}

The exponential proliferation of anthropogenic space debris within Low Earth Orbit (LEO) presents an existential hazard to orbital operations and the long-term sustainability of space exploration. Multi-decade forecasting of massive fragmentation clouds is severely bottlenecked by traditional numerical integration schemes, which suffer from cumulative truncation errors, artificial secular energy drift, and an intractable $\mathcal{O}(N^2)$ computational complexity during close-approach screening. This doctoral thesis resolves these challenges by developing, implementing, and validating an Artificial Intelligence (AI)-enhanced, high-fidelity, hybrid semi-symplectic computational framework capable of deterministic multi-decade orbital propagation and autonomous event reconstruction.

The research follows a three-stage methodological progression across three core scientific contributions. Initially, the thesis validates a scalable, structure-preserving symplectic $N$-body baseline utilizing the Wisdom-Holman Fast (WHFast) solver to propagate massive fragment clouds ($N = 20,000$) under Earth geopotential harmonics ($J_2$ and $J_4$). This baseline achieves sub-kilometer positional consistency with the National Aeronautics and Space Administration (NASA) General Mission Analysis Tool (GMAT), bounds relative energy error within $|\Delta E/E_0| \in [10^{-8}, 10^{-7}]$, and preserves an empirical power-law scalability exponent of 1.60. 

Subsequently, to incorporate non-Hamiltonian environmental dissipation without compromising symplectic invariants, the framework implements a second-order symmetric Strang operator splitting sequence. This formulation couples conservative gravitational flows with piecewise-exponential atmospheric drag and Solar Radiation Pressure (SRP) across populations exceeding 50,000 fragments, reproducing secular nodal precession within 0.005\% of analytical theory and capturing the physical ``orbital sorting'' effect driven by variations in area-to-mass ratios ($A/m$).

Finally, the research introduces an embedded AI ecosystem to eliminate computational bottlenecks in Space Traffic Management (STM). A Physics-Informed Neural Network (PINN) serves as a $\mathcal{C}^\infty$-differentiable thermospheric density surrogate, removing step-size numerical spikes caused by table lookups. Concurrently, a spatial $K$-Dimensional Tree (KD-Tree) combined with a Long Short-Term Memory (LSTM) sequence classifier prunes conjunction manifolds to $\mathcal{O}(N \log N)$ complexity. Lastly, the inverse problem of fragmentation event reconstruction is formulated as a Markov Decision Process (MDP) and solved via a Proximal Policy Optimization (PPO) agent interacting with a C-compiled symplectic engine through a shared-memory pointer bridge (\texttt{ctypes}) at 37 Frames Per Second (FPS) across 16 Central Processing Unit (CPU) cores. These results demonstrate that bridging structure-preserving astrodynamics with embedded machine learning provides an accurate, scalable tool for Space Situational Awareness (SSA) and space sustainability policy enforcement.

\vspace{0.8cm}
\noindent\textbf{Keywords:} Space Debris, Orbital Propagation, Symplectic Integrators, Strang Splitting, Physics-Informed Neural Networks (PINN), Conjunction Assessment, Reinforcement Learning, Space Traffic Management.